\documentclass[UTF8, 12pt, fontset=fandol, oneside]{ctexart}
\usepackage{researchnotes}

\title{强化学习：导论}
\author{}
\date{}

\begin{document}

\maketitle

\section*{第 5 章\quad 蒙特卡洛方法}

在本章中，我们考察第一类用于估计价值函数并发现最优策略的学习方法。
与上一章不同，这里我们不假设完全知道环境。
蒙特卡洛方法只需要\textbf{经验}，也就是来自与环境进行实际或模拟交互时得到的状态、动作和奖励的样本序列。
从\textbf{实际}经验中学习非常引人注目，因为它不需要预先知道环境动态，却仍然能够得到最优行为。
从\textbf{模拟}经验中学习同样很有力量。
虽然这需要一个模型，但模型只需要生成样本转移，而不需要给出动态规划（DP）所要求的所有可能转移的完整概率分布。
在令人惊讶的许多情形中，按照所需概率分布生成抽样经验很容易，而显式地得到这些分布却不可行。

蒙特卡洛方法是基于对样本回报取平均来求解强化学习问题的方法。
为了确保有定义良好的回报可用，这里我们只把蒙特卡洛方法定义在分幕任务上。
也就是说，我们假设经验被划分为一幕一幕，并且无论选择什么动作，所有幕最终都会终止。
只有在一幕完成之后，价值估计和策略才会被改变。
因此，蒙特卡洛方法可以在逐幕意义上是增量式的，但不能在逐步（在线）意义上是增量式的。
“蒙特卡洛”这个术语常被更宽泛地用于任何操作中含有显著随机成分的估计方法。
这里我们专门用它指基于完整回报取平均的方法，以区别于下一章讨论的、从部分回报中学习的方法。

蒙特卡洛方法为每个状态--动作对抽样并平均\textbf{回报}，这很像第 2 章中赌博机方法为每个动作抽样并平均\textbf{奖励}。
主要区别在于，现在有多个状态，每个状态像一个不同的赌博机问题（例如关联搜索或上下文赌博机），而这些不同的赌博机问题彼此关联。
也就是说，在一个状态中采取动作之后得到的回报，依赖于同一幕中后来各状态所采取的动作。
由于所有动作选择都在学习中，从较早状态的角度看，问题就变成了非平稳的。

为了处理这种非平稳性，我们采用第 4 章中为 DP 提出的广义策略迭代（GPI）思想。
在那里，我们根据 MDP 的知识\textbf{计算}价值函数；而这里，我们用来自 MDP 的样本回报来\textbf{学习}价值函数。
价值函数与相应策略仍然以本质上相同的方式相互作用，从而达到最优性（GPI）。
和 DP 章节一样，我们先考虑预测问题，即给定一个固定任意策略 $\pi$，计算 $v_\pi$ 和 $q_\pi$；然后考虑策略改进；最后考虑控制问题及其 GPI 解法。
这些来自 DP 的思想都会被扩展到只有样本经验可用的蒙特卡洛情形。

\subsection*{5.1\quad 蒙特卡洛预测}

我们先考虑用于学习给定策略下状态价值函数的蒙特卡洛方法。
回忆一下，状态的价值是从该状态开始的期望回报，也就是期望的未来累计折扣奖励。
因此，从经验中估计它的一个显然办法，就是简单地平均访问该状态之后观察到的回报。
随着观察到的回报越来越多，这个平均值应当收敛到期望价值。
这个想法正是所有蒙特卡洛方法的基础。

具体地说，假设我们想估计 $v_\pi(s)$，即在策略 $\pi$ 下状态 $s$ 的价值，并且给定了一组遵循 $\pi$、经过 $s$ 的幕。
在一幕中每一次出现状态 $s$，都称为对 $s$ 的一次\textbf{访问}。
当然，$s$ 可以在同一幕中被访问多次；我们把一幕中第一次访问它的时刻称为对 $s$ 的\textbf{首次访问}。
\textbf{首次访问 MC 方法}把首次访问 $s$ 之后的回报平均起来，作为 $v_\pi(s)$ 的估计；而\textbf{每次访问 MC 方法}则平均所有访问 $s$ 之后的回报。
这两种蒙特卡洛（MC）方法非常相似，但理论性质略有不同。
首次访问 MC 的研究最为广泛，可以追溯到 20 世纪 40 年代，本章也将重点讨论它。
每次访问 MC 更自然地扩展到函数近似和资格迹，如第 9 章和第 12 章所讨论的那样。
首次访问 MC 的过程形式如下框所示。
每次访问 MC 与它相同，只是去掉对 $S_t$ 是否已在本幕较早出现过的检查。

\begin{center}
\fbox{\begin{minipage}{0.88\textwidth}
\noindent\textbf{首次访问 MC 预测，用于估计 $V\approx v_\pi$}

\vspace{0.4em}
\noindent 输入：要评价的策略 $\pi$。

\noindent 初始化：
\[
V(s)\in\mathbb{R}\quad \text{任意，对所有 }s\in\mathcal{S},
\]
\[
Returns(s)\leftarrow \text{空列表}\quad \text{对所有 }s\in\mathcal{S}.
\]

\noindent 永远循环（对每一幕）：

\quad 按照 $\pi$ 生成一幕：
\[
S_0,A_0,R_1,S_1,A_1,R_2,\ldots,S_{T-1},A_{T-1},R_T.
\]
\quad $G\leftarrow 0$。

\quad 对本幕中的每一步 $t=T-1,T-2,\ldots,0$ 循环：

\qquad $G\leftarrow \gamma G+R_{t+1}$。

\qquad 若 $S_t$ 没有出现在 $S_0,S_1,\ldots,S_{t-1}$ 中：

\qquad\quad 将 $G$ 追加到 $Returns(S_t)$。

\qquad\quad $V(S_t)\leftarrow \mathrm{average}(Returns(S_t))$。
\end{minipage}}
\end{center}

首次访问 MC 和每次访问 MC 都会在对 $s$ 的访问次数（或首次访问次数）趋于无穷时收敛到 $v_\pi(s)$。
对首次访问 MC 来说，这一点很容易看出。
在这种情况下，每个回报都是对 $v_\pi(s)$ 的一个独立同分布估计，并且方差有限。
由大数定律，这些估计的平均值序列会收敛到其期望值。
每个平均值本身都是无偏估计，其误差的标准差随着 $1/\sqrt{n}$ 下降，其中 $n$ 是平均的回报个数。
每次访问 MC 没有这么直接，但它的估计也会二次收敛到 $v_\pi(s)$（Singh and Sutton, 1996）。

蒙特卡洛方法的使用最好通过例子来说明。

\noindent\textbf{例 5.1：二十一点}
流行的赌场纸牌游戏 blackjack（二十一点）的目标，是让手中牌面点数之和尽可能大，但不能超过 21。
所有人头牌都按 10 点计，A 可以按 1 点或 11 点计。
考虑这样一种版本：每个玩家都独立地与庄家对局。
游戏开始时，庄家和玩家各发两张牌。
庄家的一张牌正面朝上，另一张牌背面朝上。
如果玩家一开始就有 21 点（一张 A 和一张 10 点牌），这称为 natural。
此时玩家获胜，除非庄家也有 natural，在这种情况下游戏为平局。
如果玩家没有 natural，那么他可以一张一张地要牌（hit），直到他停止（stick）或超过 21 点（go bust）。
如果他爆牌，则他输；如果他停牌，则轮到庄家行动。
庄家按照一个没有选择的固定策略要牌或停牌：点数之和为 17 或更大时停牌，否则要牌。
如果庄家爆牌，则玩家获胜；否则，胜、负或平由谁的最终点数更接近 21 来决定。

玩二十一点自然可以表述为一个分幕有限 MDP。
每一局二十一点就是一幕。
获胜、失败和平局的奖励分别为 $+1$、$-1$ 和 $0$。
一局内的所有奖励都为零，并且我们不折扣（$\gamma=1$）；因此这些终止奖励也就是回报。
玩家的动作是要牌或停牌。
状态依赖于玩家手中的牌和庄家亮出的牌。
我们假设牌是从无限牌堆中发出的（即有放回抽样），因此记录已经发出的牌没有优势。
如果玩家持有一张可以按 11 点而不爆牌的 A，则称这张 A 是可用的。
在这种情况下，它总是按 11 点计，因为若按 1 点计，总和将为 11 或更小；显然此时没有决策可做，玩家总应要牌。
因此，玩家根据三个变量来决策：当前点数之和（12--21）、庄家亮出的一张牌（A--10），以及自己是否持有可用 A。
这总共给出 200 个状态。

考虑这样一个策略：若玩家点数之和为 20 或 21 则停牌，否则要牌。
为了用蒙特卡洛方法求这个策略的状态价值函数，可以用该策略模拟许多局二十一点，并平均每个状态之后的回报。
用这种方式，我们得到了图 5.1 所示的状态价值函数估计。
带有可用 A 的状态估计较不确定、也较不规则，因为这些状态较少见。
无论如何，在 500000 局之后，价值函数已经被很好地近似了。

\begin{figure}[htbp]
\centering
\includegraphics[width=0.82\textwidth]{tex/强化学习/images/math-6-fig-5-1.png}
\caption{只在 20 或 21 点停牌的二十一点策略的近似状态价值函数，它们由蒙特卡洛策略评价计算得到。}
\end{figure}

\noindent\textbf{练习 5.1}
考虑图 5.1 右侧的图。
为什么估计价值函数会在后方最后两行突然升高？
为什么左侧图中最后一整行会下降？
为什么上方图中的最前排价值高于下方图中的最前排价值？

\noindent\textbf{练习 5.2}
假设在二十一点任务中使用每次访问 MC，而不是首次访问 MC。
你预期结果会非常不同吗？
为什么，或为什么不会？

虽然在二十一点任务中我们完全知道环境，但用 DP 方法计算价值函数并不容易。
DP 方法要求下一事件的分布；特别地，它要求由四参数函数 $p$ 给出的环境动态，而这对二十一点并不容易确定。
例如，假设玩家点数之和为 14，并且他选择停牌。
作为庄家亮出牌的函数，他以奖励 $+1$ 终止的概率是多少？
在应用 DP 之前，必须先计算所有这些概率，而这样的计算通常复杂且容易出错。
相反，生成蒙特卡洛方法所需的样本牌局很容易。
这种情况出乎意料地常见；蒙特卡洛方法只凭样本幕工作的能力，即使在完全知道环境动态时，也可能是一个重要优势。

我们能否把备份图的思想推广到蒙特卡洛算法？
备份图的一般思想，是在顶部显示要更新的根节点，并在下面显示所有对该更新有贡献的转移和叶节点，其奖励和估计价值都会参与更新。
对于 $v_\pi$ 的蒙特卡洛估计，根是一个状态节点，而它下面是某个单独幕中的整条转移轨迹，最终到达终止状态，如右图所示。

\begin{figure}[htbp]
\centering
\includegraphics[height=0.28\textheight]{tex/强化学习/images/math-6-mc-backup.png}
\caption{用于蒙特卡洛估计的一条样本幕备份图。}
\end{figure}

DP 图（第 59 页）显示所有可能转移，而蒙特卡洛图只显示这一幕中实际抽样到的那些转移。
DP 图只包含一步转移，而蒙特卡洛图一直延伸到这一幕结束。
图中的这些差异准确地反映了两类算法之间的根本差异。

关于蒙特卡洛方法，一个重要事实是：每个状态的估计彼此独立。
某个状态的估计并不建立在任何其他状态的估计之上，而 DP 中则是这样。
换句话说，蒙特卡洛方法并不进行我们在上一章定义的\textbf{自举}。

特别注意，估计单个状态价值的计算代价与状态个数无关。
当人们只需要一个状态或一部分状态的价值时，这会使蒙特卡洛方法特别有吸引力。
可以从感兴趣的状态出发生成许多样本幕，只平均这些状态的回报，而忽略所有其他状态。
这是蒙特卡洛方法相对于 DP 方法可能具有的第三个优势，前两个优势是能够从实际经验和模拟经验中学习。

\noindent\textbf{例 5.2：肥皂泡}
假设有一个形成闭合回路的线框，把它浸入肥皂水中，会形成一个贴合线框边缘的肥皂膜或肥皂泡。
如果线框的几何形状不规则但已知，怎样计算这个表面的形状？
这个形状具有这样的性质：邻近点施加在每一点上的总力为零，否则形状就会改变。
这意味着表面在任一点的高度，等于该点周围一个小圆上各点高度的平均值。
此外，表面必须在边界处与线框相接。
这类问题通常的处理办法，是在表面覆盖的区域上放置一个网格，并通过迭代计算求出各网格点的高度。
边界处的网格点被强制固定在线框上，其余点则向其四个最近邻高度的平均值调整。
这个过程随后反复迭代，很像 DP 的迭代策略评价，最终收敛到所需表面的一个良好近似。

\begin{figure}[htbp]
\centering
\includegraphics[width=0.52\textwidth]{tex/强化学习/images/math-6-soap-bubble.png}
\caption{线框上的肥皂泡。}
\end{figure}

这类似于蒙特卡洛方法最初被设计来求解的那类问题。
不要使用上面描述的迭代计算，而是想象你站在这个表面上做随机游走，以相等概率从一个网格点随机迈向相邻网格点，直到到达边界。
事实证明，边界处高度的期望值，是所需表面在起点处高度的一个良好近似；实际上，它正是上面所述迭代方法计算出的值。
因此，人们可以通过从该点开始的许多随机游走的边界高度平均值，来很好地近似某一点的高度。
如果只关心一个点的值，或任何固定的小规模点集的值，那么这种蒙特卡洛方法会比基于局部一致性的迭代方法高效得多。

\subsection*{5.2\quad 动作价值的蒙特卡洛估计}

如果没有模型，那么估计\textbf{动作价值}，即状态--动作对的价值，通常比估计\textbf{状态价值}更有用。
有了模型，只有状态价值也足以确定策略；人们只需向前看一步，并选择会导致最佳奖励与下一状态组合的动作，正如我们在 DP 章节中所做的那样。
然而，没有模型时，只有状态价值并不足够。
为了让价值能够用于建议策略，必须显式估计每个动作的价值。
因此，蒙特卡洛方法的一个主要目标就是估计 $q_*$。
为了做到这一点，我们先考虑动作价值的策略评价问题。

动作价值的策略评价问题，是估计 $q_\pi(s,a)$，即从状态 $s$ 开始、采取动作 $a$、随后遵循策略 $\pi$ 时的期望回报。
用于这个问题的蒙特卡洛方法与刚才用于状态价值的方法本质上相同，只是现在讨论的是对状态--动作对的访问，而不是对状态的访问。
如果在一幕中状态 $s$ 曾被访问并且动作 $a$ 在其中被采取，那么称状态--动作对 $s,a$ 在这一幕中被访问。
每次访问 MC 方法把某个状态--动作对的价值估计为其所有访问之后回报的平均值。
首次访问 MC 方法则平均每一幕中第一次访问该状态并选择该动作之后的回报。
和以前一样，随着每个状态--动作对的访问次数趋于无穷，这些方法会二次收敛到真实期望价值。

唯一的复杂之处在于，许多状态--动作对可能永远不会被访问。
如果 $\pi$ 是确定性策略，那么遵循 $\pi$ 时，在每个状态只能观察到其中一个动作的回报。
没有回报可供平均时，其他动作的蒙特卡洛估计就不会随经验而改进。
这是一个严重问题，因为学习动作价值的目的正是帮助我们在每个状态可用的动作之间做选择。
为了比较备选动作，我们需要估计每个状态中\textbf{所有}动作的价值，而不仅仅是当前偏好的那个动作。

这就是第 2 章在 $k$ 臂赌博机问题语境中讨论过的\textbf{保持探索}的一般问题。
为了让动作价值的策略评价有效，我们必须保证持续探索。
一种办法是指定每一幕都从一个状态--动作对开始，并且每一对都有非零概率被选为起点。
这保证了在无限多幕的极限中，所有状态--动作对都会被访问无限多次。
我们把这称为\textbf{探索起始}假设。

探索起始假设有时很有用，但一般当然不能依赖它，尤其是在直接从与环境的实际交互中学习时。
在那种情况下，起始条件不太可能如此方便。
保证所有状态--动作对都会遇到的最常见替代方法，是只考虑那些在每个状态都以非零概率选择所有动作的随机策略。
我们会在后面的章节中讨论这种方法的两个重要变体。
目前，我们保留探索起始假设，并完成一个完整蒙特卡洛控制方法的介绍。

\noindent\textbf{练习 5.3}
$q_\pi$ 的蒙特卡洛估计对应的备份图是什么？

\subsection*{5.3\quad 蒙特卡洛控制}

现在我们可以考虑如何把蒙特卡洛估计用于控制，也就是用于近似最优策略。
总体思想是按照 DP 章节中的同一模式前进，也就是按照广义策略迭代（GPI）的思想前进。
在 GPI 中，人们同时维护一个近似策略和一个近似价值函数。
价值函数被反复改变，使其更接近当前策略的价值函数；策略则相对于当前价值函数被反复改进，如右图所示。
这两类改变在某种程度上彼此相互牵制，因为每一类都会为另一类制造一个移动目标；但二者合在一起会使策略和价值函数都趋近最优。

\begin{figure}[htbp]
\centering
\includegraphics[width=0.34\textwidth]{tex/强化学习/images/math-6-gpi-action.png}
\caption{基于动作价值函数的广义策略迭代示意图。}
\end{figure}

首先，我们考虑经典策略迭代的一个蒙特卡洛版本。
在这个方法中，我们交替执行完整的策略评价和策略改进步骤，从一个任意策略 $\pi_0$ 开始，最终得到最优策略和最优动作价值函数：
\[
\pi_0 \xrightarrow{\mathrm{E}} q_{\pi_0}
\xrightarrow{\mathrm{I}} \pi_1
\xrightarrow{\mathrm{E}} q_{\pi_1}
\xrightarrow{\mathrm{I}} \pi_2
\xrightarrow{\mathrm{E}} \cdots
\xrightarrow{\mathrm{I}} \pi_*
\xrightarrow{\mathrm{E}} q_*,
\]
其中 $\xrightarrow{\mathrm{E}}$ 表示一次完整的策略评价，$\xrightarrow{\mathrm{I}}$ 表示一次完整的策略改进。
策略评价完全按照上一节所述进行。
经历许多幕之后，近似动作价值函数渐近地接近真实函数。
暂且假设我们确实观察到无限多幕，并且这些幕还通过探索起始生成。
在这些假设下，蒙特卡洛方法会对任意 $\pi_k$ 精确计算每个 $q_{\pi_k}$。

策略改进通过使策略相对于当前价值函数贪心来完成。
在这里我们有的是\textbf{动作}价值函数，因此不需要模型来构造贪心策略。
对于任意动作价值函数 $q$，相应的贪心策略是这样一个策略：对每个 $s\in\mathcal{S}$，确定性地选择动作价值最大的动作：
\begin{equation}
\pi(s)\doteq \arg\max_a q(s,a).
\end{equation}
于是，可以把每个 $\pi_{k+1}$ 构造为相对于 $q_{\pi_k}$ 的贪心策略，从而完成策略改进。
策略改进定理（第 4.2 节）可应用于 $\pi_k$ 和 $\pi_{k+1}$，因为对所有 $s\in\mathcal{S}$，
\[
\begin{aligned}
q_{\pi_k}(s,\pi_{k+1}(s))
&= q_{\pi_k}\!\left(s,\arg\max_a q_{\pi_k}(s,a)\right)\\
&= \max_a q_{\pi_k}(s,a)\\
&\ge q_{\pi_k}(s,\pi_k(s))\\
&\ge v_{\pi_k}(s).
\end{aligned}
\]
正如上一章所讨论的，这个定理保证每个 $\pi_{k+1}$ 一致地优于 $\pi_k$，或者至少与 $\pi_k$ 一样好；在后一种情况下，二者都是最优策略。
这又保证了整个过程收敛到最优策略和最优价值函数。
通过这种方式，蒙特卡洛方法只凭样本幕、而不需要关于环境动态的其他知识，就能用来寻找最优策略。

为了方便得到蒙特卡洛方法的收敛保证，我们在上面作了两个不太现实的假设。
一个是假设各幕具有探索起始，另一个是假设策略评价可以用无限多幕来完成。
为了得到实际算法，我们必须去掉这两个假设。
关于第一个假设，我们推迟到本章稍后再考虑。

现在我们关注“策略评价使用无限多幕”这个假设。
这个假设相对容易去掉。
事实上，即便在经典 DP 方法中，例如迭代策略评价，也会出现同样的问题：它们也只是渐近收敛到真实价值函数。
在 DP 和蒙特卡洛两种情况下，都有两种办法解决这个问题。
一种是牢牢坚持每次策略评价都要近似 $q_{\pi_k}$ 的想法。
通过测量和假设，得到估计误差大小及其概率的上界，然后在每次策略评价中采取足够多步骤，以保证这些界足够小。
从保证收敛到某个近似水平的意义上说，这种方法或许可以做到完全令人满意。
然而，除了最小的问题之外，它也很可能需要太多幕，因而在实践中没有用。

避免策略评价名义上所需无限多幕的第二种办法，是放弃在回到策略改进之前完成策略评价的企图。
在每个评价步骤中，我们让价值函数朝 $q_{\pi_k}$ 移动，但并不指望它真的接近，除非经过许多步骤。
我们在第 4.6 节首次介绍 GPI 思想时就使用了这个想法。
这个想法的一个极端形式是价值迭代，其中每一步策略改进之间只进行一次迭代策略评价的迭代。
价值迭代的原地版本甚至更极端；在那里，我们对单个状态交替执行改进和评价步骤。

对于蒙特卡洛策略评价，自然的做法是在逐幕基础上交替进行评价和改进。
每一幕之后，观测到的回报被用于策略评价，然后在该幕中访问过的所有状态上改进策略。
沿着这些思路得到的一个完整简单算法，称为 Monte Carlo ES，即带探索起始的蒙特卡洛算法，其伪代码如下。

\begin{center}
\fbox{\begin{minipage}{0.88\textwidth}
\noindent\textbf{Monte Carlo ES（探索起始），用于估计 $\pi\approx\pi_*$}

\vspace{0.4em}
\noindent 初始化：
\[
\pi(s)\in\mathcal{A}(s)\quad \text{任意，对所有 }s\in\mathcal{S},
\]
\[
Q(s,a)\in\mathbb{R}\quad \text{任意，对所有 }s\in\mathcal{S},\ a\in\mathcal{A}(s),
\]
\[
Returns(s,a)\leftarrow \text{空列表}\quad \text{对所有 }s\in\mathcal{S},\ a\in\mathcal{A}(s).
\]

\noindent 永远循环（对每一幕）：

\quad 随机选择 $S_0\in\mathcal{S}$、$A_0\in\mathcal{A}(S_0)$，使所有状态--动作对的概率都 $>0$。

\quad 从 $S_0,A_0$ 开始并遵循 $\pi$ 生成一幕：
\[
S_0,A_0,R_1,\ldots,S_{T-1},A_{T-1},R_T.
\]
\quad $G\leftarrow 0$。

\quad 对本幕中的每一步 $t=T-1,T-2,\ldots,0$ 循环：

\qquad $G\leftarrow \gamma G+R_{t+1}$。

\qquad 若二元组 $S_t,A_t$ 没有出现在 $S_0,A_0,S_1,A_1,\ldots,S_{t-1},A_{t-1}$ 中：

\qquad\quad 将 $G$ 追加到 $Returns(S_t,A_t)$。

\qquad\quad $Q(S_t,A_t)\leftarrow \mathrm{average}(Returns(S_t,A_t))$。

\qquad\quad $\pi(S_t)\leftarrow \arg\max_a Q(S_t,a)$。
\end{minipage}}
\end{center}

\noindent\textbf{练习 5.4}
Monte Carlo ES 的伪代码是低效的，因为对每个状态--动作对，它都维护所有回报的列表，并反复计算它们的均值。
更高效的做法是使用类似第 2.4 节中说明的技术，对每个状态--动作对只维护均值和计数，并以增量方式更新它们。
请说明伪代码应如何修改才能做到这一点。

在 Monte Carlo ES 中，每个状态--动作对的所有回报都会被累计并平均，而不管观察到这些回报时正在执行什么策略。
很容易看出，Monte Carlo ES 不可能收敛到任何次优策略。
如果它收敛到了次优策略，那么价值函数最终会收敛到该策略的价值函数，而这又会导致策略改变。
只有当策略和价值函数都最优时，稳定性才会达到。
随着动作价值函数的变化随时间减小，收敛到这个最优不动点似乎不可避免，但目前尚未被形式化证明。
在我们看来，这是强化学习中最基本的开放理论问题之一（关于一个部分解，见 Tsitsiklis, 2002）。

\noindent\textbf{例 5.3：求解二十一点}
把 Monte Carlo ES 应用于二十一点是直接的。
因为这些幕全都是模拟牌局，所以很容易安排包含所有可能性的探索起始。
在这种情况下，只需以相等概率随机选择庄家的牌、玩家点数之和，以及玩家是否有可用 A。
作为初始策略，我们使用前一个二十一点例子中评价过的策略，也就是只在 20 或 21 点停牌。
初始动作价值函数可以对所有状态--动作对都设为零。
图 5.2 显示了由 Monte Carlo ES 找到的二十一点最优策略。
这个策略与 Thorp（1966）的“基本”策略相同，唯一例外是有可用 A 的策略中最左侧的凹口；Thorp 的策略中没有这个凹口。
我们不确定这种差异的原因，但确信这里所示确实是我们所描述的二十一点版本的最优策略。

\begin{figure}[htbp]
\centering
\includegraphics[width=0.78\textwidth]{tex/强化学习/images/math-6-fig-5-2.png}
\caption{由 Monte Carlo ES 得到的二十一点最优策略和状态价值函数。图中状态价值函数由 Monte Carlo ES 得到的动作价值函数计算而来。}
\end{figure}

\subsection*{5.4\quad 无探索起始的蒙特卡洛控制}

我们如何避免探索起始这个不现实的假设？
确保所有动作被选择无限多次的唯一一般方法，是让智能体持续选择它们。
有两种办法可以保证这一点，它们导致我们称为同策略方法和离策略方法的两类方法。
同策略方法试图评价或改进用于作出决策的策略，而离策略方法评价或改进一个不同于生成数据所用策略的策略。
上面发展的 Monte Carlo ES 方法是同策略方法的一个例子。
在本节中，我们说明如何设计一种不使用不现实的探索起始假设的同策略蒙特卡洛控制方法。
离策略方法将在下一节讨论。

在同策略控制方法中，策略通常是\textbf{软}的，意思是对所有 $s\in\mathcal{S}$ 和所有 $a\in\mathcal{A}(s)$ 都有 $\pi(a|s)>0$，但它会逐渐向确定性最优策略移动。
第 2 章讨论的许多方法都为此提供了机制。
本节给出的同策略方法使用 $\varepsilon$-贪心策略，意思是它们在大多数时候选择估计动作价值最大的动作，但以概率 $\varepsilon$ 随机选择一个动作。
也就是说，所有非贪心动作都被赋予最小选择概率 $\frac{\varepsilon}{|\mathcal{A}(s)|}$，剩余的大部分概率 $1-\varepsilon+\frac{\varepsilon}{|\mathcal{A}(s)|}$ 被赋予贪心动作。
$\varepsilon$-贪心策略是 $\varepsilon$-软策略的例子；$\varepsilon$-软策略定义为：对某个 $\varepsilon>0$，所有状态和动作都满足 $\pi(a|s)\ge \frac{\varepsilon}{|\mathcal{A}(s)|}$。
在 $\varepsilon$-软策略中，$\varepsilon$-贪心策略在某种意义上是最接近贪心的策略。

同策略蒙特卡洛控制的总体思想仍然是 GPI。
和 Monte Carlo ES 一样，我们使用首次访问 MC 方法来估计当前策略的动作价值函数。
然而，没有探索起始假设时，我们不能简单地通过让策略相对于当前价值函数贪心来改进策略，因为那会阻止对非贪心动作的进一步探索。
幸运的是，GPI 并不要求策略一路采取到贪心策略，只要求它朝贪心策略移动。
在我们的同策略方法中，我们只把它移动到一个 $\varepsilon$-贪心策略。
对于任意 $\varepsilon$-软策略 $\pi$，任意相对于 $q_\pi$ 的 $\varepsilon$-贪心策略都保证优于或等于 $\pi$。
完整算法如下框所示。

\begin{center}
\fbox{\begin{minipage}{0.88\textwidth}
\noindent\textbf{同策略首次访问 MC 控制（用于 $\varepsilon$-软策略），估计 $\pi\approx\pi_*$}

\vspace{0.4em}
\noindent 算法参数：小的 $\varepsilon>0$。

\noindent 初始化：

\quad $\pi\leftarrow$ 一个任意的 $\varepsilon$-软策略。

\quad $Q(s,a)\in\mathbb{R}$ 任意，对所有 $s\in\mathcal{S},\ a\in\mathcal{A}(s)$。

\quad $Returns(s,a)\leftarrow$ 空列表，对所有 $s\in\mathcal{S},\ a\in\mathcal{A}(s)$。

\noindent 永远重复（对每一幕）：

\quad 按照 $\pi$ 生成一幕：
\[
S_0,A_0,R_1,\ldots,S_{T-1},A_{T-1},R_T.
\]
\quad $G\leftarrow 0$。

\quad 对本幕中的每一步 $t=T-1,T-2,\ldots,0$ 循环：

\qquad $G\leftarrow \gamma G+R_{t+1}$。

\qquad 若二元组 $S_t,A_t$ 没有出现在 $S_0,A_0,S_1,A_1,\ldots,S_{t-1},A_{t-1}$ 中：

\qquad\quad 将 $G$ 追加到 $Returns(S_t,A_t)$。

\qquad\quad $Q(S_t,A_t)\leftarrow \mathrm{average}(Returns(S_t,A_t))$。

\qquad\quad $A^*\leftarrow \arg\max_a Q(S_t,a)$（平局任意打破）。

\qquad\quad 对所有 $a\in\mathcal{A}(S_t)$：
\[
\pi(a|S_t)\leftarrow
\begin{cases}
1-\varepsilon+\varepsilon/|\mathcal{A}(S_t)|, & a=A^*,\\
\varepsilon/|\mathcal{A}(S_t)|, & a\ne A^*.
\end{cases}
\]
\end{minipage}}
\end{center}

相对于 $q_\pi$ 的任意 $\varepsilon$-贪心策略都是对任意 $\varepsilon$-软策略 $\pi$ 的改进，这一点由策略改进定理保证。
令 $\pi'$ 为这个 $\varepsilon$-贪心策略。
策略改进定理的条件成立，因为对任意 $s\in\mathcal{S}$，
\begin{equation}
\begin{aligned}
q_\pi(s,\pi'(s))
&= \sum_a \pi'(a|s)q_\pi(s,a)\\
&= \frac{\varepsilon}{|\mathcal{A}(s)|}\sum_a q_\pi(s,a)
 +(1-\varepsilon)\max_a q_\pi(s,a)\\
&\ge
\frac{\varepsilon}{|\mathcal{A}(s)|}\sum_a q_\pi(s,a)
 +(1-\varepsilon)\sum_a
\frac{\pi(a|s)-\frac{\varepsilon}{|\mathcal{A}(s)|}}{1-\varepsilon}
 q_\pi(s,a)\\
&= \frac{\varepsilon}{|\mathcal{A}(s)|}\sum_a q_\pi(s,a)
- \frac{\varepsilon}{|\mathcal{A}(s)|}\sum_a q_\pi(s,a)
+ \sum_a \pi(a|s)q_\pi(s,a)\\
&= v_\pi(s).
\end{aligned}
\end{equation}
括号中的求和是一个加权平均，其权重非负且和为 1，因此它必然小于或等于被平均数中的最大值。
于是，由策略改进定理，$\pi'\ge \pi$，也就是对所有 $s\in\mathcal{S}$ 都有 $v_{\pi'}(s)\ge v_\pi(s)$。
我们现在证明，只有当 $\pi'$ 和 $\pi$ 都是在 $\varepsilon$-软策略中最优时，等号才可能成立；也就是说，它们优于或等于所有其他 $\varepsilon$-软策略。

考虑一个新环境，它与原环境完全相同，唯一不同是把“策略必须是 $\varepsilon$-软的”这一要求移入环境内部。
新环境具有与原环境相同的动作集和状态集，并按如下方式行为。
如果在状态 $s$ 采取动作 $a$，那么新环境以概率 $1-\varepsilon$ 完全像旧环境一样行为。
以概率 $\varepsilon$，它会以相等概率随机重新选择动作，然后用这个新的随机动作像旧环境一样行为。
在这个新环境中使用一般策略所能达到的最优效果，与在原环境中使用 $\varepsilon$-软策略所能达到的最优效果相同。
令 $\tilde v_*$ 和 $\tilde q_*$ 表示新环境的最优价值函数。
那么，当且仅当 $v_\pi=\tilde v_*$ 时，策略 $\pi$ 在 $\varepsilon$-软策略中最优。
由 $\tilde v_*$ 的定义可知，它是下式的唯一解：
\[
\begin{aligned}
\tilde v_*(s)
&= (1-\varepsilon)\max_a \tilde q_*(s,a)
 + \frac{\varepsilon}{|\mathcal{A}(s)|}\sum_a \tilde q_*(s,a)\\
&= (1-\varepsilon)\max_a \sum_{s',r}p(s',r|s,a)\left[r+\gamma\tilde v_*(s')\right]\\
&\quad + \frac{\varepsilon}{|\mathcal{A}(s)|}\sum_a\sum_{s',r}
p(s',r|s,a)\left[r+\gamma\tilde v_*(s')\right].
\end{aligned}
\]
当等号成立并且 $\varepsilon$-软策略 $\pi$ 不再被改进时，我们还由式（5.2）知道
\[
\begin{aligned}
v_\pi(s)
&= (1-\varepsilon)\max_a q_\pi(s,a)
 + \frac{\varepsilon}{|\mathcal{A}(s)|}\sum_a q_\pi(s,a)\\
&= (1-\varepsilon)\max_a \sum_{s',r}p(s',r|s,a)\left[r+\gamma v_\pi(s')\right]\\
&\quad + \frac{\varepsilon}{|\mathcal{A}(s)|}\sum_a\sum_{s',r}
p(s',r|s,a)\left[r+\gamma v_\pi(s')\right].
\end{aligned}
\]
然而，这个方程与前一个方程相同，只是把 $\tilde v_*$ 替换为 $v_\pi$。
由于 $\tilde v_*$ 是唯一解，所以必有 $v_\pi=\tilde v_*$。

本质上，我们在过去几页中已经说明，策略迭代对 $\varepsilon$-软策略也是有效的。
使用 $\varepsilon$-软策略的自然贪心概念时，每一步都保证改进，除非已经找到了 $\varepsilon$-软策略中的最优策略。
这个分析与每一阶段中动作价值函数是如何确定的无关，但它确实假设这些函数被精确计算。
这把我们带到了与上一节大致相同的位置。
现在我们只得到了 $\varepsilon$-软策略中的最佳策略；但另一方面，我们已经去掉了探索起始假设。

\subsection*{5.5\quad 通过重要性采样进行离策略预测}

所有学习控制方法都面临一个两难困境：它们试图学习在后续\textbf{最优}行为条件下的动作价值，但为了探索所有动作（以\textbf{找到}最优动作），又需要以非最优方式行动。
它们如何在按照探索性策略行动的同时，学习关于最优策略的知识？
上一节的同策略方法实际上是一种折中：它学习的不是最优策略的动作价值，而是一个仍会探索的近似最优策略的动作价值。
更直接的方法是使用两个策略：一个是被学习、并会成为最优策略的策略；另一个更具探索性，用来生成行为。
被学习的策略称为\textbf{目标策略}，用来生成行为的策略称为\textbf{行为策略}。
在这种情况下，我们说学习来自目标策略“之外”的数据，整个过程称为\textbf{离策略学习}。

在本书余下部分，我们会同时考虑同策略方法和离策略方法。
同策略方法通常更简单，因此先考虑。
离策略方法需要额外的概念和记号，并且由于数据来自不同策略，离策略方法的方差通常更大，收敛更慢。
另一方面，离策略方法更强大、更一般。
它们包含同策略方法作为目标策略和行为策略相同的特殊情形。
离策略方法在应用中还有许多额外用途。
例如，它们常可用于从传统的非学习控制器生成的数据中学习，或从人类专家的数据中学习。
也有人认为，离策略学习是学习世界动态多步预测模型的关键（见第 17.2 节；Sutton, 2009；Sutton et al., 2011）。

本节中，我们从预测问题开始研究离策略方法，其中目标策略和行为策略都是固定的。
也就是说，假设我们希望估计 $v_\pi$ 或 $q_\pi$，但手中只有遵循另一个策略 $b$ 得到的幕，其中 $b\ne\pi$。
在这种情况下，$\pi$ 是目标策略，$b$ 是行为策略，并且二者都被认为是固定且给定的。

为了用来自 $b$ 的幕估计 $\pi$ 的价值，我们要求 $\pi$ 下采取的每个动作在 $b$ 下也至少偶尔会被采取。
也就是说，我们要求 $\pi(a|s)>0$ 蕴含 $b(a|s)>0$。
这称为\textbf{覆盖}假设。
覆盖意味着，在与 $\pi$ 不同的状态中，$b$ 必须是随机的。
另一方面，目标策略 $\pi$ 可以是确定性的；事实上，在控制应用中这是一种特别有趣的情形。
在控制中，目标策略通常是相对于当前动作价值函数估计的确定性贪心策略。
当行为策略保持随机且更具探索性（例如 $\varepsilon$-贪心策略）时，这个策略会变成确定性最优策略。
不过，在本节中，我们考虑预测问题，其中 $\pi$ 不变且给定。

几乎所有离策略方法都利用\textbf{重要性采样}，这是一种通用技术，用于在给定来自另一个分布的样本时，估计一个分布下的期望值。
我们把重要性采样应用于离策略学习，做法是按照轨迹在目标策略和行为策略下发生的相对概率来加权回报，这个相对概率称为\textbf{重要性采样比}。
给定起始状态 $S_t$，在任意策略 $\pi$ 下，后续状态--动作轨迹 $A_t,S_{t+1},A_{t+1},\ldots,S_T$ 发生的概率为
\[
\begin{aligned}
&\Pr\{A_t,S_{t+1},A_{t+1},\ldots,S_T\mid S_t,A_{t:T-1}\sim\pi\}\\
&\quad = \pi(A_t|S_t)p(S_{t+1}|S_t,A_t)\pi(A_{t+1}|S_{t+1})\cdots p(S_T|S_{T-1},A_{T-1})\\
&\quad = \prod_{k=t}^{T-1}\pi(A_k|S_k)p(S_{k+1}|S_k,A_k),
\end{aligned}
\]
其中这里的 $p$ 是式（3.4）定义的状态转移概率函数。
因此，该轨迹在目标策略与行为策略下的相对概率，即重要性采样比，为
\begin{equation}
\rho_{t:T-1}\doteq
\frac{\prod_{k=t}^{T-1}\pi(A_k|S_k)p(S_{k+1}|S_k,A_k)}
{\prod_{k=t}^{T-1}b(A_k|S_k)p(S_{k+1}|S_k,A_k)}
=\prod_{k=t}^{T-1}\frac{\pi(A_k|S_k)}{b(A_k|S_k)}.
\end{equation}
虽然轨迹概率依赖于 MDP 的转移概率，而这些概率通常未知，但它们在分子和分母中完全相同，因此相互抵消。
重要性采样比最终只依赖于两个策略和序列，而不依赖于 MDP。

回忆一下，我们希望估计目标策略下的期望回报（价值），但手中只有由于行为策略而得到的回报 $G_t$。
这些回报具有错误的期望 $\mathbb{E}[G_t|S_t=s]=v_b(s)$，因此不能通过直接平均得到 $v_\pi$。
这正是重要性采样发挥作用的地方。
比值 $\rho_{t:T-1}$ 会把回报转换为具有正确期望值的量：
\begin{equation}
\mathbb{E}[\rho_{t:T-1}G_t\mid S_t=s]=v_\pi(s).
\end{equation}

现在我们可以给出一个蒙特卡洛算法：它对一批遵循策略 $b$ 的观测幕中的回报取平均，以估计 $v_\pi(s)$。
这里方便的做法，是用跨越幕边界递增的方式给时间步编号。
也就是说，如果批中的第一幕在时间 100 于终止状态结束，那么下一幕从 $t=101$ 开始。
这样我们可以用时间步编号指代特定幕中的特定步骤。
特别地，我们可以定义所有访问状态 $s$ 的时间步集合，记为 $\mathcal{T}(s)$。
这是针对每次访问方法的；对于首次访问方法，$\mathcal{T}(s)$ 只会包含在其所在幕中首次访问 $s$ 的那些时间步。
再令 $T(t)$ 表示时间 $t$ 之后的第一个终止时刻，令 $G_t$ 表示从 $t$ 到 $T(t)$ 的回报。
于是 $\{G_t\}_{t\in\mathcal{T}(s)}$ 是属于状态 $s$ 的回报，而 $\{\rho_{t:T(t)-1}\}_{t\in\mathcal{T}(s)}$ 是相应的重要性采样比。
为了估计 $v_\pi(s)$，我们只需用这些比值缩放回报并平均结果：
\begin{equation}
V(s)\doteq
\frac{\sum_{t\in\mathcal{T}(s)}\rho_{t:T(t)-1}G_t}{|\mathcal{T}(s)|}.
\end{equation}
当重要性采样以这种简单平均方式完成时，它称为\textbf{普通重要性采样}。

一个重要替代方法是\textbf{加权重要性采样}，它使用如下定义的加权平均：
\begin{equation}
V(s)\doteq
\frac{\sum_{t\in\mathcal{T}(s)}\rho_{t:T(t)-1}G_t}
{\sum_{t\in\mathcal{T}(s)}\rho_{t:T(t)-1}},
\end{equation}
若分母为零，则估计值也取零。
为了理解这两类重要性采样，考虑它们的首次访问方法在从状态 $s$ 观察到一个单独回报之后给出的估计。
在加权平均估计中，单个回报的比值 $\rho_{t:T(t)-1}$ 在分子和分母中抵消，因此估计等于观测到的回报，与该比值无关（假设比值非零）。
鉴于这是唯一观测到的回报，这是一个合理估计，但它的期望是 $v_b(s)$ 而不是 $v_\pi(s)$，因此在统计意义上它是有偏的。
相比之下，普通重要性采样估计量（5.5）的首次访问版本在期望上总是 $v_\pi(s)$，即它是无偏的，但它可能很极端。
假设比值为 10，表示观测到的轨迹在目标策略下发生的可能性是在行为策略下的 10 倍。
在这种情况下，普通重要性采样估计将是观测回报的 10 倍。
也就是说，尽管这条幕轨迹被认为非常能代表目标策略，它的估计仍会离观测回报相当远。

形式上，两种首次访问重要性采样方法之间的差异表现在偏差和方差上。
普通重要性采样无偏，而加权重要性采样有偏，尽管其偏差渐近收敛到零。
另一方面，普通重要性采样的方差一般无界，因为这些比值的方差可以无界；而在加权估计量中，任何单个回报上的最大权重都是 1。
事实上，在假设回报有界时，即使这些比值本身的方差为无穷，加权重要性采样估计量的方差也会收敛到零（Precup, Sutton, and Dasgupta 2001）。
实践中，加权估计量通常具有显著更低的方差，因此被强烈偏好。
尽管如此，我们不会完全放弃普通重要性采样，因为它更容易扩展到本书第二部分将探讨的使用函数近似的近似方法。

普通重要性采样和加权重要性采样的每次访问方法也都是有偏的；不过同样地，随着样本数增加，偏差会渐近下降到零。
实践中，每次访问方法常常更受偏好，因为它们不需要记录哪些状态已经被访问过，并且更容易扩展到近似方法。
下一节第 110 页给出了一个使用加权重要性采样进行离策略策略评价的完整每次访问 MC 算法。

\noindent\textbf{练习 5.5}
考虑一个 MDP，它有一个单一的非终止状态和一个单一动作；该动作以概率 $p$ 转移回这个非终止状态，并以概率 $1-p$ 转移到终止状态。
令所有转移上的奖励都为 $+1$，并令 $\gamma=1$。
假设你观察到一幕，它持续了 10 步，回报为 10。
这个非终止状态价值的首次访问估计量和每次访问估计量分别是多少？

\noindent\textbf{例 5.4：二十一点中一个状态价值的离策略估计}
我们同时应用普通重要性采样方法和加权重要性采样方法，用离策略数据估计一个单独二十一点状态（例 5.1）的价值。
回忆一下，蒙特卡洛方法的一个优势，是它们可以评价单个状态而不需要为任何其他状态形成估计。
在这个例子中，我们评价的状态是：庄家亮出一张 2，玩家牌面和为 13，并且玩家有一张可用 A，也就是玩家持有一张 A 和一张 2，或等价地持有三张 A。
数据通过从这个状态开始，然后以相等概率随机选择要牌或停牌来生成，这就是行为策略。
目标策略与例 5.1 一样，只在总和为 20 或 21 时停牌。
在目标策略下，这个状态的价值约为 $-0.27726$；这个数值是通过另外生成一亿幕目标策略下的牌局并平均其回报得到的。
使用随机策略生成 1000 个离策略幕之后，两种离策略方法都能较接近地近似这个值。
为了确认它们能可靠地做到这一点，我们进行了 100 次独立运行，每次都从零估计开始并学习 10000 幕。
图 5.3 显示了所得学习曲线，即每种方法的估计平方误差随幕数变化的情况，并对 100 次运行取平均。
两个算法的误差都趋近于零，但加权重要性采样方法在开始时误差低得多，这也是实践中的典型情况。

\begin{figure}[htbp]
\centering
\includegraphics[width=0.78\textwidth]{tex/强化学习/images/math-6-fig-5-3.png}
\caption{加权重要性采样在用离策略幕估计单个二十一点状态价值时产生更低误差。}
\end{figure}

\noindent\textbf{例 5.5：无限方差}
只要经过缩放的回报具有无限方差，普通重要性采样的估计通常也会具有无限方差，从而表现出令人不满意的收敛性质；这很容易发生在轨迹含有循环的离策略学习中。
图 5.4 的插图给出了一个简单例子。
其中只有一个非终止状态 $s$ 和两个动作，right 与 left。
right 动作确定性地转移到终止状态；left 动作以概率 $0.9$ 转移回 $s$，以概率 $0.1$ 转移到终止状态。
后一种转移的奖励为 $+1$，其余奖励为 0。
考虑总是选择 left 的目标策略。
该策略下的所有幕都由若干次（可能为零次）回到 $s$ 的转移组成，随后以奖励和回报 $+1$ 终止。
因此，在目标策略下 $s$ 的价值为 $1$（$\gamma=1$）。
假设我们用离策略数据估计这个价值，行为策略以相等概率选择 right 和 left。

\begin{figure}[htbp]
\centering
\includegraphics[width=0.82\textwidth]{tex/强化学习/images/math-6-fig-5-4.png}
\caption{普通重要性采样在插图所示的一状态 MDP 上产生出人意料地不稳定的估计。这里正确估计为 $1$（$\gamma=1$），并且即使这是样本回报经重要性采样后的期望值，样本方差仍为无穷，估计并不收敛到该值。结果对应离策略首次访问 MC。}
\end{figure}

图 5.4 的下半部分显示了首次访问 MC 算法使用普通重要性采样时的 10 次独立运行。
即使经过数百万幕，估计仍未能收敛到正确值 1。
相比之下，加权重要性采样算法在第一幕以 left 动作结束之后，将永远给出恰为 1 的估计。
所有不等于 1 的回报，也就是以 right 动作结束的回报，都与目标策略不一致，因此其 $\rho_{t:T(t)-1}$ 为零，并且既不贡献给式（5.6）的分子，也不贡献给分母。
加权重要性采样算法只对与目标策略一致的回报形成加权平均，而这些回报全都恰好为 1。

我们可以通过一个简单计算验证，在这个例子中，经重要性采样缩放后的回报方差为无穷。
任意随机变量 $X$ 的方差，是其偏离均值 $\bar X$ 的平方的期望，可写为
\[
\mathrm{Var}[X]\doteq
\mathbb{E}\!\left[(X-\bar X)^2\right]
=\mathbb{E}[X^2-2X\bar X+\bar X^2]
=\mathbb{E}[X^2]-\bar X^2.
\]
因此，如果均值有限，就像本例中那样，当且仅当该随机变量平方的期望为无穷时，方差为无穷。
所以我们只需证明经重要性采样缩放后的回报的期望平方为无穷：
\[
\mathbb{E}_b\!\left[
\left(
\prod_{t=0}^{T-1}\frac{\pi(A_t|S_t)}{b(A_t|S_t)}G_0
\right)^2
\right].
\]
为了计算这个期望，我们按幕长和终止方式分情况。
首先注意，任何以 right 动作结束的幕，其重要性采样比都为零，因为目标策略永远不会采取这个动作；因此这些幕对期望没有贡献（括号中的量为零），可以忽略。
我们只需考虑这样的幕：它们包含若干次（可能为零次）left 动作并转移回非终止状态，随后再由一次 left 动作转移到终止状态。
所有这些幕的回报都为 1，因此可以忽略 $G_0$ 因子。
为了得到期望平方，我们只需对每种幕长，把该幕发生的概率乘以其重要性采样比的平方，再相加：
\[
\begin{aligned}
&= \frac12\cdot 0.1\left(\frac{1}{0.5}\right)^2
\quad\text{（长度为 1 的幕）}\\
&\quad + \frac12\cdot 0.9\cdot\frac12\cdot 0.1
\left(\frac{1}{0.5}\frac{1}{0.5}\right)^2
\quad\text{（长度为 2 的幕）}\\
&\quad + \frac12\cdot 0.9\cdot\frac12\cdot 0.9\cdot\frac12\cdot 0.1
\left(\frac{1}{0.5}\frac{1}{0.5}\frac{1}{0.5}\right)^2
\quad\text{（长度为 3 的幕）}\\
&\quad + \cdots\\
&=0.1\sum_{k=0}^{\infty}0.9^k\cdot 2^k\cdot 2
=0.2\sum_{k=0}^{\infty}1.8^k
=\infty.
\end{aligned}
\]

\noindent\textbf{练习 5.6}
如果给定的是使用 $b$ 生成的回报，那么对动作价值 $Q(s,a)$ 而不是状态价值 $V(s)$，与式（5.6）对应的方程是什么？

\noindent\textbf{练习 5.7}
在图 5.3 所示的学习曲线中，误差通常会随训练而下降，普通重要性采样方法确实如此。
但对于加权重要性采样方法，误差先升高后下降。
你认为为什么会这样？

\noindent\textbf{练习 5.8}
例 5.5 和图 5.4 中的结果使用了首次访问 MC 方法。
假设在同一问题上改用每次访问 MC 方法。
估计量的方差仍然会是无穷吗？
为什么，或为什么不会？

\subsection*{5.6\quad 增量式实现}

蒙特卡洛预测方法可以在逐幕基础上增量实现，所用技术是第 2 章（第 2.4 节）所述技术的扩展。
第 2 章中我们平均的是奖励，而蒙特卡洛方法中我们平均的是回报。
在所有其他方面，与第 2 章中相同的方法都可用于同策略蒙特卡洛方法。
对于离策略蒙特卡洛方法，我们需要分别考虑使用普通重要性采样的方法和使用加权重要性采样的方法。

在普通重要性采样中，回报由重要性采样比 $\rho_{t:T(t)-1}$（5.3）缩放，然后像式（5.5）那样简单平均。
对这些方法，我们仍可使用第 2 章的增量方法，只需用缩放后的回报代替那一章中的奖励。
剩下的是使用加权重要性采样的离策略方法。
此时必须形成回报的加权平均，因此需要一个稍有不同的增量算法。

假设我们有一列回报 $G_1,G_2,\ldots,G_{n-1}$，它们全都从同一状态开始，并且每个回报都有相应的随机权重 $W_i$，例如 $W_i=\rho_{t_i:T(t_i)-1}$。
我们希望形成估计
\begin{equation}
V_n\doteq
\frac{\sum_{k=1}^{n-1}W_kG_k}{\sum_{k=1}^{n-1}W_k},
\quad n\ge 2,
\end{equation}
并在得到一个额外回报 $G_n$ 时保持它最新。
除了跟踪 $V_n$，我们还必须为每个状态维护给前 $n$ 个回报的权重累计和 $C_n$。
$V_n$ 的更新规则为
\begin{equation}
V_{n+1}\doteq V_n+\frac{W_n}{C_n}\left[G_n-V_n\right],
\quad n\ge 1,
\end{equation}
并且
\[
C_{n+1}\doteq C_n+W_{n+1},
\]
其中 $C_0=0$，而 $V_1$ 任意，因此无需指定。
下一页的框给出了一个完整的逐幕增量算法，用于蒙特卡洛策略评价。
该算法名义上用于离策略情形，使用加权重要性采样；但它同样适用于同策略情形，只需把目标策略和行为策略选为相同策略即可，此时 $\pi=b$，$W$ 总为 1。
在动作按照可能不同的策略 $b$ 选择时，近似 $Q$ 会对所有遇到的状态--动作对收敛到 $q_\pi$。

\noindent\textbf{练习 5.9}
修改第 5.1 节中首次访问 MC 策略评价的算法，使其使用第 2.4 节所述样本平均的增量式实现。

\noindent\textbf{练习 5.10}
从式（5.7）推导加权平均更新规则（5.8）。
请遵循非加权规则（2.3）的推导模式。

\begin{center}
\fbox{\begin{minipage}{0.88\textwidth}
\noindent\textbf{离策略 MC 预测（策略评价），用于估计 $Q\approx q_\pi$}

\vspace{0.4em}
\noindent 输入：任意目标策略 $\pi$。

\noindent 初始化，对所有 $s\in\mathcal{S},\ a\in\mathcal{A}(s)$：
\[
Q(s,a)\in\mathbb{R}\quad\text{任意},\qquad C(s,a)\leftarrow 0.
\]

\noindent 永远循环（对每一幕）：

\quad $b\leftarrow$ 任意覆盖 $\pi$ 的策略。

\quad 按照 $b$ 生成一幕：
\[
S_0,A_0,R_1,\ldots,S_{T-1},A_{T-1},R_T.
\]
\quad $G\leftarrow 0$，$W\leftarrow 1$。

\quad 对本幕中的每一步 $t=T-1,T-2,\ldots,0$ 循环，同时要求 $W\ne 0$：

\qquad $G\leftarrow \gamma G+R_{t+1}$。

\qquad $C(S_t,A_t)\leftarrow C(S_t,A_t)+W$。

\qquad $Q(S_t,A_t)\leftarrow Q(S_t,A_t)+
\dfrac{W}{C(S_t,A_t)}\left[G-Q(S_t,A_t)\right]$。

\qquad $W\leftarrow W\dfrac{\pi(A_t|S_t)}{b(A_t|S_t)}$。
\end{minipage}}
\end{center}

\subsection*{5.7\quad 离策略蒙特卡洛控制}

现在我们可以展示本书考虑的第二类学习控制方法的一个例子：离策略方法。
回忆一下，同策略方法的区别性特征，是它们在使用某个策略进行控制的同时估计该策略的价值。
在离策略方法中，这两个功能被分离。
用于生成行为的策略称为行为策略，它事实上可以与被评价和改进的策略无关；后者称为目标策略。
这种分离的一个优势是，目标策略可以是确定性的（例如贪心的），而行为策略可以持续抽样所有可能动作。

离策略蒙特卡洛控制方法使用前两节介绍的技术之一。
它们遵循行为策略，同时学习并改进目标策略。
这些技术要求行为策略以非零概率选择目标策略可能选择的所有动作，即满足覆盖。
为了探索所有可能性，我们要求行为策略是软的，也就是说，它在所有状态中都以非零概率选择所有动作。

下一页的框给出了一种离策略蒙特卡洛控制方法，它基于 GPI 和加权重要性采样，用于估计 $\pi_*$ 和 $q_*$。
目标策略 $\pi\approx\pi_*$ 是相对于 $Q$ 的贪心策略，而 $Q$ 是对 $q_\pi$ 的估计。
行为策略 $b$ 可以是任意策略，但为了保证 $\pi$ 收敛到最优策略，必须为每个状态--动作对得到无限多个回报。
这可以通过选择 $\varepsilon$-软的 $b$ 来保证。
即使动作是按照一个不同的软策略 $b$ 选择的，而且 $b$ 可能在幕与幕之间甚至幕内发生变化，策略 $\pi$ 也会在所有遇到的状态上收敛到最优。

\begin{center}
\fbox{\begin{minipage}{0.88\textwidth}
\noindent\textbf{离策略 MC 控制，用于估计 $\pi\approx\pi_*$}

\vspace{0.4em}
\noindent 初始化，对所有 $s\in\mathcal{S},\ a\in\mathcal{A}(s)$：
\[
Q(s,a)\in\mathbb{R}\quad\text{任意},\qquad C(s,a)\leftarrow 0,
\]
\[
\pi(s)\leftarrow \arg\max_a Q(s,a)\quad\text{（平局一致地打破）}.
\]

\noindent 永远循环（对每一幕）：

\quad $b\leftarrow$ 任意软策略。

\quad 使用 $b$ 生成一幕：
\[
S_0,A_0,R_1,\ldots,S_{T-1},A_{T-1},R_T.
\]
\quad $G\leftarrow 0$，$W\leftarrow 1$。

\quad 对本幕中的每一步 $t=T-1,T-2,\ldots,0$ 循环：

\qquad $G\leftarrow \gamma G+R_{t+1}$。

\qquad $C(S_t,A_t)\leftarrow C(S_t,A_t)+W$。

\qquad $Q(S_t,A_t)\leftarrow Q(S_t,A_t)+
\dfrac{W}{C(S_t,A_t)}\left[G-Q(S_t,A_t)\right]$。

\qquad $\pi(S_t)\leftarrow \arg\max_a Q(S_t,a)$（平局一致地打破）。

\qquad 如果 $A_t\ne \pi(S_t)$，则退出内层循环（进入下一幕）。

\qquad $W\leftarrow W\dfrac{1}{b(A_t|S_t)}$。
\end{minipage}}
\end{center}

一个潜在问题是，这个方法只从幕的尾部学习，也就是从该幕剩余动作全都是贪心动作的位置开始学习。
如果非贪心动作很常见，那么学习会很慢，特别是对出现在长幕早期部分的状态而言。
这可能会大大拖慢学习。
目前关于离策略蒙特卡洛方法的经验还不足以评估这个问题有多严重。
如果它确实严重，那么解决它的最重要办法可能是引入时序差分学习，也就是下一章发展的算法思想。
另一种可能是，如果 $\gamma<1$，那么下一节提出的思想也可能有显著帮助。

\noindent\textbf{练习 5.11}
在离策略 MC 控制的框中算法里，你可能原本预期 $W$ 的更新会涉及重要性采样比 $\frac{\pi(A_t|S_t)}{b(A_t|S_t)}$，但它实际涉及的是 $\frac{1}{b(A_t|S_t)}$。
为什么这仍然是正确的？

\noindent\textbf{练习 5.12：赛车场（编程）}
考虑驾驶赛车通过类似图 5.5 所示的弯道。
你希望尽可能快地通过，但又不能快到冲出赛道。
在我们的简化赛车场中，赛车位于一个离散网格位置集合中的某一格，即图中的格子。
速度也是离散的，即每个时间步在水平方向和竖直方向移动的网格数。
动作是速度分量的增量。
每个速度分量在每一步都可以改变 $+1$、$-1$ 或 $0$，因此共有九个（$3\times3$）动作。
两个速度分量都被限制为非负且小于 5，并且除起跑线外二者不能同时为零。
每一幕从随机选择的某个起始状态开始，此时两个速度分量都为零；当赛车越过终点线时，幕结束。
赛车越过终点线之前，每一步奖励为 $-1$。
如果赛车撞到赛道边界，它会被移回起跑线上的一个随机位置，两个速度分量都被置零，并且这一幕继续。
在每个时间步更新赛车位置之前，检查赛车的投影路径是否与赛道边界相交。
如果它与终点线相交，则幕结束；如果它在其他地方相交，则认为赛车撞到赛道边界并被送回起跑线。
为了使任务更具挑战性，在每个时间步，以概率 $0.1$，两个速度增量都会变为零，与原本打算的增量无关。
请把一种蒙特卡洛控制方法应用到这个任务，以计算从每个起始状态出发的最优策略。
展示若干条遵循最优策略的轨迹，但这些轨迹中请关闭噪声。

\begin{figure}[htbp]
\centering
\includegraphics[width=0.78\textwidth]{tex/强化学习/images/math-6-fig-5-5.png}
\caption{赛车场任务中的两个右转弯。}
\end{figure}

\subsection*{5.8\quad *折扣感知的重要性采样}

到目前为止，我们考虑的离策略方法，都是对被视为整体的回报形成重要性采样权重，并没有考虑回报作为折扣奖励之和的内部结构。
现在我们简要考虑一种前沿研究思想：利用这种结构显著降低离策略估计量的方差。

例如，考虑幕很长且 $\gamma$ 显著小于 1 的情形。
为具体起见，假设幕持续 100 步且 $\gamma=0$。
那么从时间 0 开始的回报只是 $G_0=R_1$，但它的重要性采样比却是 100 个因子的乘积：
\[
\frac{\pi(A_0|S_0)}{b(A_0|S_0)}
\frac{\pi(A_1|S_1)}{b(A_1|S_1)}
\cdots
\frac{\pi(A_{99}|S_{99})}{b(A_{99}|S_{99})}.
\]
在普通重要性采样中，回报会由整个乘积来缩放，但其实只需要用第一个因子 $\frac{\pi(A_0|S_0)}{b(A_0|S_0)}$ 来缩放。
其他 99 个因子与回报无关，因为第一个奖励之后回报已经确定。
这些后续因子都与回报独立，并且期望值为 1；它们不会改变期望更新，却会极大增加其方差。
在某些情况下，它们甚至可能使方差变为无穷。
现在我们考虑一个避免这种巨大额外方差的想法。

这个想法的本质，是把折扣看作决定终止概率，或者等价地，决定部分终止程度。
对任意 $\gamma\in[0,1)$，我们可以认为回报 $G_0$ 以程度 $1-\gamma$ 在一步后部分终止，产生只含第一个奖励 $R_1$ 的回报；又以程度 $(1-\gamma)\gamma$ 在两步后部分终止，产生 $R_1+R_2$ 的回报，依此类推。
后一个程度对应于第二步继续的程度 $1-\gamma$，并且已经在第一步继续的程度为 $\gamma$。
因此，第三步终止的程度为 $(1-\gamma)\gamma^2$，其中 $\gamma^2$ 反映前两步都没有终止。
这里的部分回报称为\textbf{平坦部分回报}：
\[
\bar G_{t:h}\doteq R_{t+1}+R_{t+2}+\cdots+R_h,
\quad 0\le t<h\le T,
\]
其中“平坦”表示没有折扣，“部分”表示这些回报并不一直延伸到终止，而是在称为视界的 $h$ 处停止；$T$ 是该幕的终止时刻。
传统的完整回报 $G_t$ 可以按如下方式看成平坦部分回报的和：
\[
\begin{aligned}
G_t &\doteq R_{t+1}+\gamma R_{t+2}+\gamma^2R_{t+3}+\cdots+\gamma^{T-t-1}R_T\\
&= (1-\gamma)R_{t+1}
 +(1-\gamma)\gamma(R_{t+1}+R_{t+2})\\
&\quad +(1-\gamma)\gamma^2(R_{t+1}+R_{t+2}+R_{t+3})+\cdots\\
&\quad +(1-\gamma)\gamma^{T-t-2}(R_{t+1}+R_{t+2}+\cdots+R_{T-1})\\
&\quad +\gamma^{T-t-1}(R_{t+1}+R_{t+2}+\cdots+R_T)\\
&=(1-\gamma)\sum_{h=t+1}^{T-1}\gamma^{h-t-1}\bar G_{t:h}
 +\gamma^{T-t-1}\bar G_{t:T}.
\end{aligned}
\]
现在，我们需要用同样被截断的重要性采样比来缩放这些平坦部分回报。
由于 $\bar G_{t:h}$ 只涉及到视界 $h$ 之前的奖励，所以只需要直到 $h$ 的概率比。
我们定义一个类似式（5.5）的普通重要性采样估计量：
\begin{equation}
V(s)\doteq
\frac{
\sum_{t\in\mathcal{T}(s)}
\left((1-\gamma)\sum_{h=t+1}^{T(t)-1}\gamma^{h-t-1}\rho_{t:h-1}\bar G_{t:h}
+\gamma^{T(t)-t-1}\rho_{t:T(t)-1}\bar G_{t:T(t)}\right)}
{|\mathcal{T}(s)|},
\end{equation}
以及一个类似式（5.6）的加权重要性采样估计量：
\begin{equation}
V(s)\doteq
\frac{
\sum_{t\in\mathcal{T}(s)}
\left((1-\gamma)\sum_{h=t+1}^{T(t)-1}\gamma^{h-t-1}\rho_{t:h-1}\bar G_{t:h}
+\gamma^{T(t)-t-1}\rho_{t:T(t)-1}\bar G_{t:T(t)}\right)}
{\sum_{t\in\mathcal{T}(s)}
\left((1-\gamma)\sum_{h=t+1}^{T(t)-1}\gamma^{h-t-1}\rho_{t:h-1}
+\gamma^{T(t)-t-1}\rho_{t:T(t)-1}\right)}.
\end{equation}
我们把这两个估计量称为\textbf{折扣感知的重要性采样估计量}。
它们考虑了折扣率；但若 $\gamma=1$，它们没有影响，与第 5.5 节的离策略估计量相同。

\subsection*{5.9\quad *每决策重要性采样}

还有另一种方式，可以在离策略重要性采样中考虑回报作为奖励和的结构；这种方式即使在没有折扣时（即 $\gamma=1$）也可能降低方差。
在离策略估计量（5.5）和（5.6）中，分子中求和的每一项本身都是一个和：
\begin{equation}
\begin{aligned}
\rho_{t:T-1}G_t
&=\rho_{t:T-1}\left(R_{t+1}+\gamma R_{t+2}+\cdots+\gamma^{T-t-1}R_T\right)\\
&=\rho_{t:T-1}R_{t+1}
+\gamma\rho_{t:T-1}R_{t+2}
+\cdots+\gamma^{T-t-1}\rho_{t:T-1}R_T.
\end{aligned}
\end{equation}
离策略估计量依赖于这些项的期望，而这些期望可以用更简单的方式写出。
注意，式（5.11）的每个子项都是一个随机奖励与一个随机重要性采样比的乘积。
例如，使用式（5.3），第一个子项可以写为
\begin{equation}
\rho_{t:T-1}R_{t+1}
=\frac{\pi(A_t|S_t)}{b(A_t|S_t)}
\frac{\pi(A_{t+1}|S_{t+1})}{b(A_{t+1}|S_{t+1})}
\frac{\pi(A_{t+2}|S_{t+2})}{b(A_{t+2}|S_{t+2})}
\cdots
\frac{\pi(A_{T-1}|S_{T-1})}{b(A_{T-1}|S_{T-1})}
R_{t+1}.
\end{equation}
在所有这些因子中，人们可能会怀疑只有第一个因子和最后的奖励彼此相关；所有其他因子都对应奖励之后发生的事件。
而且，所有这些其他因子的期望值都是 1：
\begin{equation}
\mathbb{E}\!\left[\frac{\pi(A_k|S_k)}{b(A_k|S_k)}\right]
\doteq \sum_a b(a|S_k)\frac{\pi(a|S_k)}{b(a|S_k)}
=\sum_a \pi(a|S_k)=1.
\end{equation}
再经过几个步骤，可以证明，正如所怀疑的那样，所有这些其他因子在期望上都没有影响，也就是说
\begin{equation}
\mathbb{E}[\rho_{t:T-1}R_{t+1}]
=\mathbb{E}[\rho_{t:t}R_{t+1}].
\end{equation}
如果对式（5.11）的第 $k$ 个子项重复这个过程，就得到
\[
\mathbb{E}[\rho_{t:T-1}R_{t+k}]
=\mathbb{E}[\rho_{t:t+k-1}R_{t+k}].
\]
于是，原始项（5.11）的期望可写为
\[
\mathbb{E}[\rho_{t:T-1}G_t]=\mathbb{E}[\tilde G_t],
\]
其中
\[
\tilde G_t=\rho_{t:t}R_{t+1}
+\gamma\rho_{t:t+1}R_{t+2}
+\gamma^2\rho_{t:t+2}R_{t+3}
+\cdots+\gamma^{T-t-1}\rho_{t:T-1}R_T.
\]
我们把这个思想称为\textbf{每决策重要性采样}。
由此立刻得到一个替代的重要性采样估计量：它在首次访问情形下与普通重要性采样估计量（5.5）具有相同的无偏期望，但使用 $\tilde G_t$：
\begin{equation}
V(s)\doteq \frac{\sum_{t\in\mathcal{T}(s)}\tilde G_t}{|\mathcal{T}(s)|},
\end{equation}
我们可以预期它有时具有更低的方差。

是否存在每决策版本的加权重要性采样？
这一点不那么清楚。
到目前为止，据我们所知，为此提出的所有估计量都不是一致的，也就是说，即使数据无穷，它们也不收敛到真实值。

\noindent\textbf{练习 5.13}
请给出从式（5.12）推导式（5.14）的步骤。

\noindent\textbf{练习 5.14}
修改第 111 页离策略蒙特卡洛控制算法，使其使用截断加权平均估计量（5.10）的思想。
注意，你首先需要把这个方程转换为动作价值形式。

\subsection*{5.10\quad 总结}

本章介绍的蒙特卡洛方法，从样本幕形式的经验中学习价值函数和最优策略。
这使它们相对于 DP 方法至少有三类优势。
第一，它们可以直接从与环境的交互中学习最优行为，而不需要环境动态模型。
第二，它们可以与模拟或样本模型一起使用。
对令人惊讶的许多应用来说，模拟样本幕很容易，而构造 DP 方法所需的转移概率显式模型则很困难。
第三，把蒙特卡洛方法聚焦在一小部分状态上既容易又高效。
一个特别感兴趣的区域可以被准确评价，而不必花代价准确评价状态集的其余部分；我们会在第 8 章进一步探讨这一点。

蒙特卡洛方法的第四个优势会在本书后面讨论：它们可能较少受到马尔可夫性质违背的伤害。
这是因为它们并不基于后继状态的价值估计来更新自己的价值估计。
换句话说，这是因为它们不自举。

在设计蒙特卡洛控制方法时，我们遵循了第 4 章介绍的广义策略迭代（GPI）总体框架。
GPI 包含相互作用的策略评价过程和策略改进过程。
蒙特卡洛方法提供了一种替代的策略评价过程。
它们不是使用模型来计算每个状态的价值，而是简单地平均从该状态开始的许多回报。
因为状态价值就是期望回报，所以这个平均值可以成为价值的良好近似。
在控制方法中，我们尤其关心近似动作价值函数，因为它们可用于改进策略，而不需要环境转移动态模型。
蒙特卡洛方法在逐幕基础上交织策略评价和策略改进，并且可以在逐幕基础上增量实现。

在蒙特卡洛控制方法中，保持\textbf{充分探索}是一个问题。
仅仅选择当前估计最好的动作是不够的，因为这样就得不到其他备选动作的回报，也可能永远学不到它们其实更好。
一种方法是假设各幕从随机选择的状态--动作对开始，以覆盖所有可能性，从而忽略这个问题。
这种\textbf{探索起始}在使用模拟幕的应用中有时可以安排，但在从真实经验中学习时不太可能。
在\textbf{同策略}方法中，智能体承诺总是探索，并试图找到仍然会探索的最佳策略。
在\textbf{离策略}方法中，智能体也会探索，但学习的是一个可能与所遵循策略无关的确定性最优策略。

\textbf{离策略预测}指的是：从另一个不同的\textbf{行为策略}生成的数据中，学习某个\textbf{目标策略}的价值函数。
这类学习方法基于某种形式的\textbf{重要性采样}，也就是根据在两个策略下采取观测动作的概率之比来加权回报，从而把期望从行为策略转换到目标策略。
\textbf{普通重要性采样}使用加权回报的简单平均，而\textbf{加权重要性采样}使用加权平均。
普通重要性采样产生无偏估计，但方差更大，且可能为无穷；加权重要性采样总是具有有限方差，在实践中更受偏好。
尽管概念上很简单，用于预测和控制的离策略蒙特卡洛方法仍未完全定论，并且仍是持续研究的主题。

本章讨论的蒙特卡洛方法与上一章讨论的 DP 方法有两个主要区别。
第一，它们作用于样本经验，因此可以用于无模型的直接学习。
第二，它们不自举。
也就是说，它们并不基于其他价值估计来更新自己的价值估计。
这两个差异并非紧密相连，而是可以分离。
下一章中，我们将考虑像蒙特卡洛方法一样从经验中学习、但又像 DP 方法一样自举的方法。

\section*{书目与历史注释}

“Monte Carlo”这个术语可以追溯到 20 世纪 40 年代，当时洛斯阿拉莫斯的物理学家设计了可研究的随机游戏，以帮助理解与原子弹相关的复杂物理现象。
关于这种意义上的蒙特卡洛方法，可参见若干教材（例如 Kalos and Whitlock, 1986；Rubinstein, 1981）。

\noindent\textbf{5.1--2}
Singh and Sutton（1996）区分了每次访问 MC 方法和首次访问 MC 方法，并证明了这些方法与强化学习算法相关的结果。
二十一点例子基于 Widrow, Gupta, and Maitra（1973）使用的一个例子。
肥皂泡例子是一个经典 Dirichlet 问题，其蒙特卡洛解首先由 Kakutani（1945）提出；见 Hersh and Griego（1969）、Doyle and Snell（1984）。

Barto and Duff（1994）在用于求解线性方程组的经典蒙特卡洛算法语境中讨论了策略评价。
他们使用 Curtiss（1954）的分析指出了蒙特卡洛策略评价在大规模问题上的计算优势。

\noindent\textbf{5.3--4}
Monte Carlo ES 在本书 1998 年版中被引入。
这可能是蒙特卡洛估计与基于策略迭代的控制方法之间第一次明确连接。
在强化学习语境中，早期使用蒙特卡洛方法估计动作价值的是 Michie and Chambers（1968）。
在杆平衡问题（第 56 页）中，他们使用幕持续时间的平均值来评估每个状态中每个可能动作的价值，即期望平衡“寿命”，然后使用这些评估来控制动作选择。
他们的方法在精神上类似于使用每次访问 MC 估计的 Monte Carlo ES。
Narendra and Wheeler（1986）研究了一种用于遍历有限马尔可夫链的蒙特卡洛方法，该方法把连续访问同一状态之间累计的回报用作奖励，以调整学习自动机的动作概率。

\noindent\textbf{5.5}
高效的离策略学习已经被认为是一个出现在多个领域中的重要挑战。
例如，它与概率图（贝叶斯）模型中的“干预”和“反事实”思想密切相关（例如 Pearl, 1995；Balke and Pearl, 1994）。
使用重要性采样的离策略方法有很长历史，但仍未被充分理解。
加权重要性采样有时也称为归一化重要性采样（例如 Koller and Friedman, 2009），Rubinstein（1981）、Hesterberg（1988）、Shelton（2001）以及 Liu（2001）等都讨论过它。

离策略学习中的目标策略在文献中有时被称为“估计”策略，本书第一版中也这样称呼。

\noindent\textbf{5.7}
赛车场练习改编自 Barto, Bradtke, and Singh（1995）以及 Gardner（1973）。

\noindent\textbf{5.8}
我们对折扣感知重要性采样思想的处理，基于 Sutton, Mahmood, Precup, and van Hasselt（2014）的分析。
迄今为止，Mahmood（2017；Mahmood, van Hasselt, and Sutton, 2014）对它给出了最完整的展开。

\noindent\textbf{5.9}
每决策重要性采样由 Precup, Sutton, and Singh（2000）引入。
他们还把离策略学习与时序差分学习、资格迹和近似方法结合起来，引入了我们将在后续章节中考虑的微妙问题。

\end{document}
